\documentclass[11pt,letterpaper]{article}
\usepackage[margin=0.72in]{geometry}
\usepackage[T1]{fontenc}
\usepackage{lmodern,microtype}
\usepackage[hidelinks]{hyperref}
\setlength{\parindent}{0pt}
\setlength{\parskip}{8pt}
\pagestyle{empty}
\begin{document}
\begin{flushright}
Rochester, New York\\
\href{mailto:garciapiterz@gmail.com}{garciapiterz@gmail.com}\\
+1 (646) 288-3300\\[12pt]
August 18, 2026
\end{flushright}

Neuromatch Hiring Team

Dear Neuromatch Hiring Team,

I am applying for the Curriculum Specialist, NeuroAI for High School position.
The role closely matches the work I am building across K--12 computer science
education, applied AI research, and inclusive curriculum design. I am completing
M.S. degrees in Teaching Computer Science K--12 at the University of Rochester
and Data Science at Rochester Institute of Technology, where I also serve as a
graduate assistant in quantum and AI research.

In my K--5 computer science placement, middle-school science teaching, and the
Freedom Scholars Planet Protectors environmental-science camp, I have learned to
turn complex material into clear goals, modeled first steps, visual and verbal
supports, structured collaboration, and multiple ways for students to demonstrate
understanding. I use backward design and Universal Design for Learning without
lowering intellectual expectations. My professional data and AI experience also
helps me build reproducible activities, documentation, and assessments that make
technical ideas inspectable rather than mysterious.

Neuromatch's plan to develop open, modular NeuroAI materials for grades 11--12 is
especially compelling because it joins scientific rigor with broad access. I
would contribute practical knowledge of machine learning, responsible AI,
curriculum design, disability-informed access, and collaboration with educators.
I am comfortable translating topics such as neural representation, feedback,
generalization, failure, fairness, and decision-making into inquiry-driven
learning while questioning the social assumptions embedded in technical systems.

The remote, approximately 15-hour weekly structure fits my current teaching and
graduate commitments, and I can participate in regular meetings during Eastern
U.S. business hours. I would be glad to help Neuromatch and its partners create,
test, document, and revise materials that teachers can adapt across classrooms.

Sincerely,

Piter Garcia
\end{document}
